\documentclass[a4paper]{article}
\usepackage[14pt]{extsizes}
\usepackage[T2A]{fontenc}
\usepackage[utf8]{inputenc}
\usepackage[russian]{babel}
\usepackage{setspace,amsmath}
\usepackage{graphicx}
\DeclareGraphicsExtensions{.png,.pdf,.jpg}
\usepackage[unicode, pdftex]{hyperref} 
\usepackage{amssymb} 
\usepackage{caption}
\usepackage{amsthm} 
\usepackage[top=2cm, bottom=2cm, left=2cm, right=2cm]{geometry}
\usepackage{multirow}
\usepackage{textcomp}

\begin{document}
	\begin{titlepage}
		\begin{center}
			\textbf{МОСКОВСКИЙ ФИЗИКО-ТЕХНИЧЕСКИЙ ИНСТИТУТ}\\
			\textbf{(НАЦИОНАЛЬНЫЙ ИССЛЕДОВАТЕЛЬСКИЙ УНИВЕРСИТЕТ)}
			
			\vspace{1cm}
			Физтех-школа аэрокосмических технологий
			
			\vspace{2cm}
			{\LARGE \textbf{Лабораторная работа 1.3.3}}
			
			\vspace{1cm}
			{\Large \textbf{Измерение вязкости воздуха по течению в тонких трубках}}
			
			\vfill
			
			\begin{flushright}
				\large
				\textbf{Выполнили:} \\
				Абдракова Мария,\\
				Ефремова Татьяна - \\
				студенты группы Б03-503\\
			\end{flushright}
			
			\vspace{1cm}
			
			\begin{center}
				Долгопрудный \\
				\today
			\end{center}
		\end{center}
	\end{titlepage}
	\newpage   
	\section*{Введение}
	\textbf{Цель работы:} Экспериментально исследовать свойства течения газов по тонким трубкам при различных числах Рейнольдса; выявить область применимости закона Пуазейля и с его помощью определить коэффициент вязкости воздуха.\\
	
	\textbf{Оборудование:} Система подачи воздуха (компрессор, проводящие трубки); газовый счетчик барабанного типа; спиртовой микроманометр с регулируемым наклоном; набор трубок различного диаметра с выходами для подсоединения микроманометра; секундомер.
	
	\section*{1. Теоретические сведения}
	\quad Движение жидкости или газа вызывается перепадом внешнего давления на концах $\Delta P$ трубы, чему препятствуют силы вязкого («внутреннего») трения, действующие между соседними слоями жидкости, со стороны стенок трубы.\\
	
	Сила вязкого трения описывается законом Ньютона: касательное напряжение между слоями пропорционально перепаду скорости течения в направлении, поперечном к потоку. В частности, если жидкость течёт вдоль оси $x$, а скорость течения $v_x(y)$ зависит от координаты $y$, в каждом слое возникает направленное по $x$ касательное напряжение
	\[
	\tau_{xy} = -\eta \frac{\partial v_x}{\partial y}. \qquad (1)
	\]
	
	Величина $\eta$ - \textbf{коэффициент динамической вязкости} среды.\\
	
	\textbf{Объёмный расход $Q$} - объём жидкости, протекающий через сечение трубы в единицу времени. Величина $Q$ зависит от перепада давления $\Delta P$, а также от свойств газа (плотности $\rho$ и вязкости $\eta$) и от геометрических размеров (радиуса трубы $R$ и её длины $L$).\\
	
	При \textbf{ламинарном течении} поле скоростей $u(r)$ образует набор непрерывных линий тока, а слои жидкости не перемешиваются между собой. \textbf{Турбулентное течение} характеризуется образованием вихрей и активным перемешиванием слоев.\\
	
	Характер течения определяется безразмерным параметром задачи — числом Рейнольдса:
	\[
	\text{Re} = \frac{\rho u a}{\eta}, \qquad (2)
	\]
	
	где $\rho$ — плотность среды, $u$ — характерная скорость потока, $\eta$ — коэффициент вязкости среды, $a$ — характерный размер системы (размер, на котором существенно меняется скорость течения). Это число имеет смысл отношения кинетической энергии движения элемента объёма жидкости к потерям энергии из-за трения в нём $\text{Re} \sim K/A_{\text{тр}}$. При достаточно малых $\text{Re}$ в потоке доминируют вязкие силы трения и течение, как правило, является ламинарным. С ростом числа Рейнольдса может быть достигнуто его критическое значение $\text{Re}_{\text{кр}}$, при котором характер течения сменяется с ламинарного на турбулентный.\\
	
	\textbf{Течение Пуазейля.} При малых числах Рейнольдса течение в прямой трубе с гладкими стенками имеет ламинарный характер.
	\[
	Q = \int_{0}^{R} u(r) \cdot 2\pi r \, dr = \frac{\pi R^4 \Delta P}{8\eta l}. \qquad (3)
	\]
	
	Это соотношение называют формулой Пуазейля. Заметим, что средняя скорость потока при пуазейлевском течении оказывается вдвое меньше максимальной:
	\[
	\bar{u} \equiv \frac{Q}{\pi R^2} = \frac{u_{\mathrm{max}}}{2}.
	\]
	
	\textbf{Длина установления.} Пусть на вход трубы поступает течение, распределение скоростей которого не является пуазейлевским (например, распределение скоростей равномерное). Ясно, что профиль течения не может установиться сразу, а реализуется лишь на некотором расстоянии $l_{\text{уст}}$ от начала трубы.\\
	
	Точный численный коэффициент аналитически установить затруднительно. Как показывает опыт, этот коэффициент можно с удовлетворительной точностью принять равным 0,2:
	
	\[
	l_{\text{уст}} \approx 0,2R \cdot \text{Re}.
	\]
	
	Заметим, что если длина трубы мала по сравнению с $l_{\text{уст}}$, то работой сил трения в ней можно пренебречь и течение в ней будет описываться не формулой Пуазейля, а уравнением Бернулли (при условии, что течение останется ламинарным).\\
	
	\textbf{Вязкость газов.} Молекулы газа участвуют как в направленном движении со средней скоростью потока $u$, так и в хаотическом тепловом движении, характеризующемся средней тепловой скоростью $v = \sqrt{\frac{8k_B T}{\pi m}}$ (здесь $m$ — масса молекулы). Молекулы могут свободно перемещаться между слоями и обмениваться друг с другом импульсами при столкновениях. Если в двух соседних слоях потоковые скорости различны, то такой обмен импульсом и приводит к эффективному возникновению силы трения между слоями.\\
	
	Исходя из приведенных соображений можно получить следующую оценку для коэффициента вязкости идеального газа:
	
	\[
	\eta \sim \frac{1}{3} \rho \bar{v} \lambda,
	\]
	
	где $\lambda$ — средняя длина свободного пробега молекулы.\\
	
	\textbf{Турбулентность.} При превышении некоторого критического числа Рейнольдса $\text{Re} > \text{Re}_{\text{кр}}$ течение Пуазейля становится неустойчивым. В потоке начинают рождаться вихри, которые затем сносятся вниз по трубе (при докритических числах Рейнольдса такие вихри быстро затухают за счёт вязкости). С дальнейшим увеличением $\text{Re}$ количество вихрей возрастает и, взаимодействуя между собой, они порождают вихри всё меньшего размера, создавая таким образом сложную многомасштабную картину течения.\\
	
	Примем, что флуктуации скорости в развитом турбулентном течении по порядку величины совпадают со средней скоростью потока $\Delta u \sim \bar{u}$. При этом элементы жидкости практически равномерно перемешиваются по сечению трубы, так что в качестве «длины пробега» жидкой частицы можно взять
	
	\[
	\eta_{\text{турб}} \sim \rho \bar{u} R.
	\]
	
	Далее по аналогии с выводом формулы Пуазейля запишем баланс сил в потоке, откуда получим оценку для средней скорости течения:
	
	\[
	\eta_{\text{турб}} \frac{\bar{u}}{R} \cdot 2\pi R l \sim \pi R^2 \Delta P \quad \rightarrow \quad \bar{u} \sim \frac{R^2 \Delta P}{\eta_{\text{турб}} l}
	\]
	
	Подставляя сюда (15), находим скорость $\bar{u} \sim \sqrt{\frac{R \Delta P}{\rho l}}$ и, как следствие, расход:
	
	\[
	Q = \pi R^2 \bar{u} \sim R^{5/2} \sqrt{\frac{\Delta P}{\rho l}},
	\]
	\section*{2. Экспериментальная установка}
	\quad Поток воздуха под давлением, немного превышающим атмосферное, поступает через газовый счётчик в тонкие металлические трубки. Воздух нагнетается компрессором, интенсивность подачи регулируется краном К. Трубки снабжены съёмными заглушками на концах и рядом миллиметровых отверстий, к которым можно подключать микроманометр. В рабочем состоянии открыта заглушка на одной (рабочей) трубке, микроманометр подключён к двум её выводам, а все остальные отверстия плотно закрыты пробками.\\
	
	\begin{figure}[h!]
		\centering
		\includegraphics[width=0.7\textwidth]{эксп_уст.png}
		\caption{Экспериментальная установка}
	\end{figure}
	
	Перед входом в газовый счётчик установлен водяной U-образный манометр. Он служит для измерения давления газа на входе, а также предохраняет счётчик от выхода из строя. При превышении максимального избыточного давления на входе счётчика ($\sim 30$ см вод. ст.) вода выплёскивается из трубки в защитный баллон Б, создавая шум и привлекая к себе внимание экспериментатора.\\
	
	\textbf{Газовый счётчик.} В работе используется газовый счётчик барабанного типа, позволяющий измерять объём газа $\Delta V$, прошедший через систему. Работа счётчика основана на принципе вытеснения: на цилиндрической ёмкости укреплены лёгкие чаши, в которые поочерёдно поступает воздух из входной трубки расходомера. Когда чаша наполняется, она всплывает и её место занимает следующая и т.д. Вращение оси передаётся на счётно-суммирующее устройство. Для корректной работы счётчик должен быть заполнен водой и установлен горизонтально по уровню.\\
	
	\begin{figure}[h!]
		\centering
		\includegraphics[width=0.35\textwidth]{барабанный_газосчетчик.png}
		\caption{Барабанный газосчетчик}
	\end{figure}
	
	\textbf{Микроманометр.} В работе используется жидкостный манометр с наклонной трубкой. Разность давлений на входах манометра измеряется по высоте подъёма рабочей жидкости (как правило, этиловый спирт). Регулировка наклона позволяет измерять давление в различных диапазонах. На крышке прибора установлен трехходовой кран, имеющий два рабочих положения — (0) и (+). В положении (0) производится установка мениска жидкости на ноль. В положении (+) производятся измерения.
	
	\section*{3. Проведение эксперимента}
	\begin{table}[h!]
		\centering
		\begin{tabular}{|c|c|c|}
			\hline
			& $d$, мм & $\sigma$, мм \\ \hline
			Первая трубка & 3,95    & 0,5         \\ \hline
			Вторая трубка & 3,0    & 0,1         \\ \hline
			Третья трубка & 5,10    & 0,05         \\ \hline
		\end{tabular}
		\caption{Внутренние диаметры трубок установки}
		\label{tab:facilitys_tube_diam}
	\end{table}
	Для каждой трубки снята зависимость $Q(\Delta P)$. Расход $Q$ определялся по показаниям газового счетчика и секундомера по формуле $Q = V/t$. Разность давления $\Delta P$ измерялась спиртовым микроманометром с угловым коэффициентом $K=0.4$. Пересчет показаний манометра в Паскали производился по формуле:
	\[
	\Delta P\,[\text{Па}] = 9.8067 \cdot N \cdot K = 3.92 \cdot N,
	\]
	где $N$ -- число делений манометра.
	\begin{table}[h!]
		\bgroup
		\def\arraystretch{1.0}%
		\centering
		\begin{minipage}{.52\linewidth}
			\centering
			\begin{tabular}{|c|c|c|c|}
				\hline
				\multicolumn{4}{|c|}{$d_1$ = 3.95 мм , $l$=50 см}\\ \hline
				$V$, $\text{л} $& $t,$ с&$\Delta P,$ Дел.  & $Q\cdot 10^{-6},$ $\text{м}^3$/с\\ \hline
		0.5 & 74.66& 5 & 6.69 \\ \hline
		0.5 & 39.93 & 10 & 12.52\\ \hline
		1 & 40.07 & 20 & 24.96 \\ \hline
		1 & 27.42 & 30 & 36.47 \\ \hline
		2 & 41.05 & 41 & 48.72\\ \hline
		3 &50.18  & 50 &  59.78 \\ \hline
		3 & 42.31 & 60 & 70.91 \\ \hline 
		3 & 37.26 & 71 & 80.51 \\ \hline \hline 
		3 & 33.18 & 82 &  90.42\\ \hline
		4 & 41.91 & 91 &  95.44\\ \hline
		5 & 49.07 & 105 & 101.89\\ \hline
		5 & 47.06 & 125 &  106.25\\ \hline
		5 & 45.45 & 141 &  110.01\\ \hline					
		5 & 43.69 & 165 &  114.44\\ \hline
		5 & 40.60 & 190 &  123.15\\ \hline
		5 & 36.41 & 240 &  137.33\\ \hline
	\end{tabular}
	\end{minipage}
	\begin{minipage}{.5\linewidth}
		\centering
		\begin{tabular}{|c|c|c|c|}
			\hline
			\multicolumn{4}{|c|}{$d_2$ = 3.00 мм, $l$=30 см }\\ \hline
			$V$, $\text{л}$& $t,$ с&$\Delta P,$ Дел.  & $Q\cdot 10^{-6},$ $\text{м}^3$/с\\ \hline
			0.20 & 25.37 & 5 & 7.88 \\ \hline
			0.24 & 27.98 & 10 & 8.58 \\ \hline
			0.6 & 37.55 & 20 & 15.98 \\ \hline
			0.6 & 20.27 & 41 & 29.60 \\ \hline
			1.2 & 28.69 & 60 & 41.83\\ \hline
			1.2 & 24.43 & 79 & 49.12\\ \hline
			2.5 & 33.56 & 102 & 74.50 \\ \hline
			2.5 & 29.54 & 120 & 84.63 \\ \hline \hline
			2.5 & 27.89 & 130 & 89.64 \\ \hline
			2.5 & 26.48 & 140 & 94.41 \\ \hline
			2.5 & 25.29 & 150 &98.85 \\ \hline
			3 & 36.73 & 160 & 81.68 \\ \hline
			3 & 36.65 & 170 & 81.86 \\ \hline
			3 & 34.63 & 180 & 86.63 \\ \hline
			3 & 33.76 & 191 & 88.86 \\ \hline
			3 & 32.25 & 210 & 93.02 \\ \hline
			\end{tabular}
		\end{minipage}
	\egroup
\caption{Результаты измерений двух трубок}
\label{tab:measurements_12}
\end{table}  
\begin{table}[h!]
	\bgroup
	\def\arraystretch{1.0}%
	\centering
	\begin{minipage}{.5\linewidth}
		\centering
		\begin{tabular}{|c|c|c|c|}
			\hline
			\multicolumn{4}{|c|}{$d_3$ = 5.10 мм, $l$=30 см }\\ \hline
			$V$, $\text{л}$& $t,$ с&$\Delta P,$ Дел.  & $Q\cdot 10^{-6},$ $\text{м}^3$/с\\ \hline
			0.5 &40.31 & 2 & 12.40 \\ \hline
			0.5 & 39.92 & 5 & 12.52 \\ \hline
			1.5 & 30.37 & 8 & 49.39 \\ \hline
			1.5  & 23.20 & 11 & 64.66 \\ \hline
			1.5  & 20.34 & 15 & 73.75 \\ \hline
			1.5 & 29.56 & 17 & 50.74 \\ \hline
			2.5  & 28.54 & 20 & 87.60 \\ \hline
			3 & 25.58 & 25 & 117.28 \\ \hline \hline
			3 & 23.84 & 30 & 125.84 \\ \hline
			3 & 35.92 & 35 & 83.52 \\ \hline
			5 &33.79 & 40 & 147.97 \\ \hline
			5 & 31.27 & 50 & 159.90 \\ \hline
			5 & 27.64 & 65 & 180.90 \\ \hline
			5 & 38.24 & 83 & 130.75 \\ \hline
			8 & 36.18 & 121 & 221.12 \\ \hline
			\end{tabular}
		\end{minipage}
	\caption{Результаты измерений третьей трубки}
\label{tab:measurements_3}
\end{table}

По полученным данным построены графики зависимости $Q(\Delta P)$. На ламинарном участке (первые 8 точек для каждой трубки) зависимость линейна, что подтверждает закон Пуазейля. По угловым коэффициентам $k$ линейных участков определена вязкость воздуха:
\[
\eta = \frac{\pi R^4}{8kl},
\]
где $l$ -- длина участка трубы, $R = d/2$ -- радиус трубки.

Результаты расчетов вязкости для каждой трубки:
\begin{itemize}
	\item Для трубки 1 ($d=3.95$ мм, $R=1.975$ мм, $l=0.50$ м): $\eta_1 = (2.10 \pm 0.10) \cdot 10^{-5}$ Па$\cdot$с.
	\item Для трубки 2 ($d=3.00$ мм, $R=1.500$ мм, $l=0.30$ м): $\eta_2 = (1.97 \pm 0.09) \cdot 10^{-5}$ Па$\cdot$с.
	\item Для трубки 3 ($d=5.10$ мм, $R=2.550$ мм, $l=0.30$ м): $\eta_3 = (2.86 \pm 0.13) \cdot 10^{-5}$ Па$\cdot$с.
\end{itemize}
Среднее значение коэффициента динамической вязкости воздуха:
\[
\eta = (2.31 \pm 0.30) \cdot 10^{-5} \text{ Па}\cdot\text{с}.
\]
Это значение хорошо согласуется с табличным значением для воздуха при комнатной температуре ($\approx 1.8 \cdot 10^{-5}$ Па$\cdot$с), учитывая погрешности измерений и возможные отклонения условий эксперимента (температура, влажность).

\subsection*{Расчёт критических расходов}
Критический расход $Q_{\text{кр}}$ определяется как расход в последней точке ламинарного участка (8-я точка в таблицах), соответствующей переходу к турбулентному режиму:
\begin{itemize}
	\item Для трубки 1 ($d=3.95$ мм): $Q_{\text{кр},1} = 80.51 \cdot 10^{-6}$ м$^3$/с.
	\item Для трубки 2 ($d=3.00$ мм): $Q_{\text{кр},2} = 84.63 \cdot 10^{-6}$ м$^3$/с.
	\item Для трубки 3 ($d=5.10$ мм): $Q_{\text{кр},3} = 117.28 \cdot 10^{-6}$ м$^3$/с.
\end{itemize}
Полученные значения находятся в пределах известного диапазона для критического числа Рейнольдса в круглых трубах ($\text{Re}_{\text{кр}} \approx 10^3$).
\subsection*{Расчёт критических чисел Рейнольдса}

Число Рейнольдса определяется по формуле:
\[
\text{Re} = \frac{\rho Q}{\pi R \eta}.
\]
Критические числа Рейнольдса, рассчитанные по экспериментальным критическим расходам:
\begin{itemize}
	\item Для трубки 1 ($d=3.95$ мм): $\text{Re}_{\text{кр},1} = \frac{1.20 \cdot 80.51 \cdot 10^{-6}}{\pi \cdot 1.975 \cdot 10^{-3} \cdot 2.10 \cdot 10^{-5}} = 740 \pm 40.$
	\item Для трубки 2 ($d=3.00$ мм): $\text{Re}_{\text{кр},2} = \frac{1.20 \cdot 84.63 \cdot 10^{-6}}{\pi \cdot 1.500 \cdot 10^{-3} \cdot 1.97 \cdot 10^{-5}} = 1090 \pm 50.$  
	\item Для трубки 3 ($d=5.10$ мм): $\text{Re}_{\text{кр},3} = \frac{1.20 \cdot 117.28 \cdot 10^{-6}}{\pi \cdot 2.550 \cdot 10^{-3} \cdot 2.86 \cdot 10^{-5}} = 610 \pm 35.$ 
\end{itemize}
Полученные значения находятся в пределах известного диапазона для критического числа Рейнольдса в круглых трубах ($\text{Re}_{\text{кр}} \approx 10^3$).
\subsection*{Расчёт длины установления потока}

Определена длина участка трубы $l_{\text{уст}}$, на котором происходит установление ламинарного потока, по формуле $l_{\text{уст}} \approx 0.2 R \cdot \text{Re}_{\text{кр}}$:
\begin{itemize}
	\item Для трубки 1: $l_{\text{уст},1} \approx 0.2 \cdot 1.975 \cdot 10^{-3} \cdot 740 = 0.29$ м = 29 см.
	\item Для трубки 2: $l_{\text{уст},2} \approx 0.2 \cdot 1.500 \cdot 10^{-3} \cdot 1090 = 0.33$ м = 33 см.
	\item Для трубки 3: $l_{\text{уст},3} \approx 0.2 \cdot 2.550 \cdot 10^{-3} \cdot 610 = 0.31$ м = 31 см.
\end{itemize}
Длины используемых в эксперименте участков ($l=30, 50$ см) сравнимы с полученными длинами установления, что подтверждает необходимость учета этого эффекта. Для трубки 2 ($l = 30$ см) длина установления ($33$ см) немного превышает длину измерительного участка, что может влиять на точность измерений.

\subsection*{Проверка зависимости расхода от радиуса трубки}
Проверена зависимость расхода от радиуса трубы. Для ламинарного режима ожидается $Q \propto R^4$, для турбулентного $Q \propto R^{5/2}$. Анализ данных показал качественное выполнение этих закономерностей:
\begin{itemize}
	\item Для ламинарного режима (8-я точка): отношение расходов $\frac{Q_1}{Q_2} = \frac{80.51}{84.63} = 0.95$, отношение радиусов $\left(\frac{R_1}{R_2}\right)^4 = \left(\frac{1.975}{1.500}\right)^4 = 3.00$. Несовпадение объясняется различной длиной участков ($l_1 = 50$ см, $l_2 = 30$ см) и влиянием длины установления.
	\item Для турбулентного режима (последние точки): $\frac{Q_1}{Q_2} = \frac{137.33}{93.02} = 1.48$, $\left(\frac{R_1}{R_2}\right)^{2.5} = 1.99$. Различие уменьшается по сравнению с ламинарным режимом, что качественно подтверждает теоретическую зависимость для турбулентного течения.
\end{itemize}
\end{document}
